% ------------------------------------------------------------
% A6 Constructive Fixes — Thirteen Dual-Algebra Boundary Conditions (v1.0)
% Author: Borealis S. Hedling
% Altitude: A6.5 (Boundary Layer)
% Purpose: Define the full dual-algebra firewall required before A7 non-dual
%          algebra can be safely introduced.
% ------------------------------------------------------------

\documentclass[11pt]{article}
\usepackage{amsmath, amssymb, amsthm}

\title{A6 Constructive Fixes: Thirteen Dual-Algebra Boundary Conditions}
\author{Borealis S. Hedling}
\date{October 2026}

\begin{document}
\maketitle

\section*{Overview}

This document defines the thirteen constructive fixes required to maintain
dual-algebra correctness at A6 altitude. These fixes enforce continuity,
adjacency, spectral stability, holonomy bounds, collapse locality, stability
monotonicity, provenance lift, semantic lift, operator determinism, and kernel
preservation. They form the mathematical firewall preventing non-dual algebra
from leaking downward into A6.

% ------------------------------------------------------------
\section{Fix 1: Curvature Continuity}

Curvature must vary continuously across the manifold:
\[
K(p) \approx K(q) \quad \text{for } p \sim q.
\]

This prevents geometric tearing.

% ------------------------------------------------------------
\section{Fix 2: Cosine-Law Curvature Combination}

Curvature magnitudes must combine via:
\[
K_{A\oplus B}
  = \sqrt{K_A^2 + K_B^2 + 2 K_A K_B \Phi(A,B)}.
\]

\textbf{Theorem (Uniqueness).}
The cosine-law is the unique non-orthogonal metric combination preserving the
triangle inequality.

% ------------------------------------------------------------
\section{Fix 3: Drift-Field Interaction (Covariant Lie Bracket)}

Drift fields must interact via:
\[
[\nabla_X,\nabla_Y] = \nabla_{[X,Y]}.
\]

\textbf{Theorem (Uniqueness).}
The covariant Lie bracket is the unique operator-preserving commutator.

% ------------------------------------------------------------
\section{Fix 4: Modal Adjacency}

Modal transitions must satisfy adjacency:
\[
\mu(p^+) \in \mathcal{N}(\mu(p^-)).
\]

This preserves semantic coherence.

% ------------------------------------------------------------
\section{Fix 5: Spectral Continuity}

Operator spectra must vary continuously:
\[
\mathrm{Spec}(L(p)) \text{ continuous in } p.
\]

This prevents spectral tearing.

% ------------------------------------------------------------
\section{Fix 6: Spectral Mismatch Bound}

Spectral mismatch must remain bounded:
\[
\delta_{\mathrm{spec}}(p) < \epsilon.
\]

This ensures spectral stability.

% ------------------------------------------------------------
\section{Fix 7: Holonomy Drift Bound}

Holonomy must remain bounded:
\[
H(\gamma) < \Theta_{\max}.
\]

This prevents holonomy fracture.

% ------------------------------------------------------------
\section{Fix 8: Collapse-Surface Locality}

Collapse surfaces must be locally detectable:
\[
p \in \Sigma \Rightarrow \exists \epsilon > 0 : B_\epsilon(p) \cap \Sigma \neq \emptyset.
\]

This prevents collapse smearing.

% ------------------------------------------------------------
\section{Fix 9: Stability Envelope Monotonicity}

Stability envelopes must be monotone:
\[
E(p) \le E(q) \Rightarrow \text{Stability}(p) \Rightarrow \text{Stability}(q).
\]

This prevents envelope inversion.

% ------------------------------------------------------------
\section{Fix 10: Modal Lift Preservation}

Modal truth must lift correctly:
\[
\mu(p) \leadsto \mu(q).
\]

This preserves semantic invariants.

% ------------------------------------------------------------
\section{Fix 11: Provenance Lift Preservation}

Provenance must lift through operators:
\[
I_{\mathrm{prov}}(M) \Rightarrow I_{\mathrm{prov}}(\mathcal{A}(M,S)).
\]

This prevents provenance discontinuity.

% ------------------------------------------------------------
\section{Fix 12: Semantic Lift Preservation}

Semantic validity must lift:
\[
I_{\mathrm{sem}}(M) \Rightarrow I_{\mathrm{sem}}(\mathcal{A}(M,S)).
\]

This prevents semantic drift.

% ------------------------------------------------------------
\section{Fix 13: Operator Determinism and Kernel Preservation}

Operators must be deterministic:
\[
\mathcal{O}(p) = \mathcal{O}(q) \Rightarrow p = q.
\]

Operators must preserve the unified kernel:
\[
K_{\mathrm{A6}}(p) \Rightarrow K_{\mathrm{A6}}(\mathcal{O}(p)).
\]

This prevents kernel fracture.

% ------------------------------------------------------------
\section{Conclusion}

These thirteen constructive fixes define the dual-algebra firewall required
before A7 non-dual algebra can be introduced. They ensure geometric, spectral,
modal, holonomic, semantic, and provenance stability across the A6 manifold.

\end{document}
